\documentclass{beamer}
\usepackage[T2A]{fontenc} % Без этого с компилятором pdflatex не работает
\usepackage[utf8]{inputenc} % Без этого с компилятором pdflatex не работает
\usepackage[russian]{babel} % Без этого с компилятором pdflatex не работает
\usetheme{Madrid} 
\usecolortheme{seahorse}
\setbeamertemplate{footline}{%
	\leavevmode%
	\hbox{%
		\begin{beamercolorbox}[wd=.6\paperwidth,ht=2.5ex,dp=1ex,left]{palette primary}%
			\hspace{0.5em}Студенты АлтГТУ, группа 8 ПИ-51
		\end{beamercolorbox}%
		\begin{beamercolorbox}[wd=.4\paperwidth,ht=2.5ex,dp=1ex,right]{palette primary}%
			\insertshortdate\hspace{1em}\insertframenumber/\inserttotalframenumber\hspace{0.5em}
		\end{beamercolorbox}%
	}%
	\vskip 0pt%
}

\title{Обзор статьи}
\subtitle{Characterization of Large Language Model Development in the Datacenter\\
	Характеризация разработки больших языковых моделей в дата-центре}

\author{
	Студенты АлтГТУ, группа 8 ПИ-51\\
	Бельш А.Е.\\
	Воробьев К.В.\\
	Мэн Ч.\\
	Ошурков Н.А.
}

\date{2026 г.}

\begin{document}
	
	
	\begin{frame}{Ключевые выводы статьи}
		\begin{itemize}
			\item LLM-нагрузки отличаются от классического DL
			\item Сеть и межузловые коммуникации становятся критическим фактором
			\item Отказы оказывают сильное влияние при длительном pretraining
			\item Существующие кластеры не оптимизированы под масштаб LLM
		\end{itemize}
		
		\vspace{0.3cm}
		\textbf{Главный результат:} обучение LLM требует специализированной
		архитектуры и системной оптимизации.
	\end{frame}
	
	\begin{frame}{Направления развития}
		\begin{itemize}
			\item Оптимизация long-sequence pretraining
			\item Улучшение поддержки MoE
			\item Повышение эффективности RLHF
			\item Модернизация сетевой инфраструктуры (NIC)
			\item Расширение кластера для более масштабных моделей
		\end{itemize}
		
		\vspace{0.3cm}
		\textbf{Итог:} статья формирует инженерные принципы построения
		дата-центров для LLM.
	\end{frame}
	
	\frame{\titlepage}
\end{document}


